\documentclass[11pt]{article}

\usepackage[margin=1.1in]{geometry}
\usepackage{amsmath,amssymb,amsthm,mathtools}
\usepackage{enumitem}
\usepackage{hyperref}
\usepackage[nameinlink,noabbrev]{cleveref}

\hypersetup{colorlinks=true,linkcolor=blue,citecolor=blue,urlcolor=blue,pdftitle={Lower bounds for the sparse travelling salesman constant of Aldous and Krikun},pdfauthor={Yitzchak Shmalo}}

\newtheorem{theorem}{Theorem}[section]
\newtheorem{lemma}[theorem]{Lemma}
\newtheorem{proposition}[theorem]{Proposition}
\newtheorem{corollary}[theorem]{Corollary}
\theoremstyle{remark}
\newtheorem{remark}[theorem]{Remark}

\providecommand{\doi}[1]{\href{https://doi.org/#1}{doi:#1}}
\newcommand{\R}{\mathbb{R}}
\newcommand{\E}{\mathbb{E}}
\newcommand{\Prob}{\mathbb{P}}

\title{Lower bounds for the sparse travelling salesman constant of Aldous and Krikun}
\author{Yitzchak Shmalo\\
Einstein Institute of Mathematics, The Hebrew University of Jerusalem\\
Givat Ram, Jerusalem, Israel}
\date{}

\begin{document}
\maketitle

\begin{abstract}
Let $T_m$ be the length of the shortest path from the origin through $m$ distinct points of a Poisson process of unit
intensity in $\R^d$, $d\ge 2$. Aldous and Krikun proved that $T_m/m$ converges in probability to a constant $c(0+)>0$ and
remarked that reasonable bounds on its value seem difficult to obtain. Their branching random walk argument, with its
integral evaluated exactly, gives $c(0+)\ge b_d=(d/e)(v_d\,d!)^{-1/d}$, where $v_d$ is the volume of the unit ball, and
$b_2=0.29352\dots$. The expected number of paths behind this bound tends to infinity for every larger threshold, and
with a theorem of Rhee $c(0+)/\sqrt d\to(2\pi e)^{-1/2}$. A shortest path cannot be shortened by
exchanging two consecutive points. Counting only the paths with this property, and bounding their number by a transfer
operator and a positive vector checked in exact arithmetic, we prove $c(0+)>0.32417$ for $d=2$.
\end{abstract}

\noindent\textbf{Keywords:} Euclidean travelling salesman problem, Poisson process, first moment method, transfer
operator, computer-assisted proof.\\
\textbf{MSC 2020:} 60D05, 60K35, 90C27.

\section{Introduction}\label{sec:intro}

Let $\Pi=(\xi_i)$ be a Poisson process of unit intensity in $\R^d$, $d\ge 2$, let
$v_d=\pi^{d/2}/\Gamma(d/2+1)$ be the volume of the unit ball in $\R^d$, where $\Gamma$ is Euler's gamma function, and for
$m\ge 1$ let
\[
T_m=\min\Bigl\{|\xi_{j_1}|+\sum_{i=2}^{m}|\xi_{j_i}-\xi_{j_{i-1}}|\Bigr\},
\]
the minimum over ordered choices of $m$ distinct points of $\Pi$. Aldous and Krikun \cite{AK} studied the smallest
average edge length of long paths through the points of $\Pi$. For $0<\delta\le 1$ let $L_k(\delta)$ be the length of the shortest cycle through
$\lceil\delta k\rceil$ of $k$ independent uniform points in a cube of volume $k$. By \cite[Theorem~1]{AK}:
\begin{enumerate}[label=(\alph*)]
\item $L_k(\delta)/(\delta k)\to c(\delta)$ in $L^2$, for a constant $c(\delta)$;
\item $c(\cdot)$ is non-decreasing and continuous on $(0,1]$, and $c(0+)=\lim_{\delta\downarrow 0}c(\delta)>0$;
\item[(e)] $T_m/m\to c(0+)$ in probability.
\end{enumerate}
We call $c(0+)$ the sparse travelling salesman constant: by (a) and (b) it is the limit, as $\delta\downarrow 0$, of
the constants for shortest cycles through a proportion $\delta$ of the points.

The constant $c(1)$ is the Beardwood--Halton--Hammersley constant $\beta_d$ of the travelling salesman problem
\cite{BHH}, so (b) gives $c(0+)\le\beta_d$. In the plane, $\beta_2\le 0.90367$ by a computer-aided argument of
Gaudio and Guan \cite{GG}. Aldous gave a construction showing $c(0+)\le 0.607$ for $d=2$, rigorous up to the sampling
error of a Monte Carlo estimate \cite{Aapp}. Open Problem~3 of \cite{AK} proposes to study the value of $c(0+)$ in
$d=2$ by Monte Carlo methods, and \cite[Section~2.1]{AK} remarks that reasonable bounds on it seem difficult to obtain.
Carlsson, Behroozi, Devulapalli and Meng \cite{CBDM} study a generalized travelling salesman problem, in which a tour
must visit one point from each of several classes of uniform random points, and compare it with the problem of
\cite{AK}. The constant $c(0+)$ also appears in the continuous greedy paths model of Gou\'er\'e and Marchand \cite{GM};
Verges \cite{Verges} studies concentration and large deviations for that model. With every point of
$\Pi$ given mass one, let $W(\ell)$ be the largest number of points of $\Pi$ on a path of length at most $\ell$ starting
at the origin. Then $W(\ell)\ge m$ exactly when $T_m\le\ell$, so the limit of $W(\ell)/\ell$ is $1/c(0+)$, and upper
bounds for that limit are lower bounds for $c(0+)$.

The known lower bounds for $c(0+)$ come from a first moment. For $a>0$ let $Z_m(a)$ be the number of ordered
$m$-tuples of distinct points $x_1,\dots,x_m$ of $\Pi$ with $\sum_{i=1}^m|x_i-x_{i-1}|\le am$, where $x_0=0$, so that
$\Prob(T_m\le am)\le\E Z_m(a)$. Aldous and Krikun \cite[Section~3.3]{AK} compare $T_m$ with a branching random walk,
and their argument bounds $\E Z_m(a)$ by
\[
\Bigl(e^{\lambda a}\int_{\R^d}e^{-\lambda|u|}\,du\Bigr)^m,\qquad \lambda>0 .
\]
The integral equals $v_d\,d!\,\lambda^{-d}$, and the choice $\lambda=d/a$ gives $c(0+)\ge b_d$, where
\[
b_d=\frac de\,(v_d\,d!)^{-1/d},\qquad b_2=\frac{\sqrt{2/\pi}}{e}=0.29352\dots,\qquad b_3=0.37677\dots .
\]
In \cite{AK} the integral is printed as $v_d\,\Gamma(d+1)\lambda^{-(d-1)}$, which leads to the value
$(2\pi e)^{-1}\approx 0.0585$ in the plane, and a later note of Aldous \cite{Anote} replaces that value by
$e(d-1)(v_d\,\Gamma(d+1))^{-1/(d-1)}$, which is $e/(2\pi)\approx 0.4326$ in the plane. By \Cref{cor:asym} this
expression exceeds $\beta_d\ge c(0+)$ for all large $d$, so it is not a lower bound for $c(0+)$ in high dimension.
Whether $c(0+)\ge e/(2\pi)$ holds in the plane is not decided by the bounds below. The proof of \cite[Lemma~2.1]{GM}
contains the same count with $\lambda=1/a$, which gives $c(0+)\ge(e\,v_d\,d!)^{-1/d}$, equal to
$(2\pi e)^{-1/2}\approx 0.2420$ in the plane. The count $Z_m(a)$ gives nothing beyond $b_d$:

\begin{theorem}\label{thm:first}
For every $d\ge 2$, $a>0$ and $m\ge 1$,
\[
\E Z_m(a)=\frac{(v_d\,d!)^m(am)^{dm}}{(dm)!}.
\]
Consequently, as $m\to\infty$, $\E Z_m(a)\to 0$ if $a<b_d$ and $\E Z_m(a)\to\infty$ if $a>b_d$; and $c(0+)\ge b_d$.
\end{theorem}

For large $d$ the lower bound $b_d$ has the same first order asymptotics as the upper bound $\beta_d$.

\begin{corollary}\label{cor:asym}
$c(0+)/\sqrt d\to(2\pi e)^{-1/2}$ as $d\to\infty$. Moreover
there is $d_0$ such that $e(d-1)(v_d\,\Gamma(d+1))^{-1/(d-1)}>\beta_d\ge c(0+)$ for all $d\ge d_0$.
\end{corollary}

The upper half of the limit is Rhee's theorem $\beta_d/\sqrt d\to(2\pi e)^{-1/2}$ \cite{Rhee}; the lower half is
$c(0+)\ge b_d$.

In the plane we pass $b_2$ by counting fewer paths. If $(0=x_0,x_1,\dots,x_m)$ is a path of length $T_m$ and
$u_i=x_i-x_{i-1}$, then exchanging $x_i$ and $x_{i+1}$ does not shorten it (this is the 2-opt move of Croes
\cite{Croes} on a segment of two points), so
\[
|u_i|+|u_{i+2}|\le|u_i+u_{i+1}|+|u_{i+1}+u_{i+2}|,\qquad 1\le i\le m-2
\]
(\Cref{lem:swap}). Let $q_m$ be the probability that $m$ independent vectors with density $e^{-|u|}/(2\pi)$ on $\R^2$
satisfy these $m-2$ inequalities. Then for every $\lambda>0$
\[
\Prob(T_m\le am)\le\bigl(e^{\lambda a}\,2\pi\lambda^{-2}\bigr)^m q_m ,
\]
and $\limsup_m q_m^{1/m}\le\rho$ implies $c(0+)\ge b_2\,\rho^{-1/2}$ (\Cref{prop:block}). The inequalities for
disjoint triples of steps are independent and each fails with probability at least $1/27$, which gives
$c(0+)\ge b_2(27/26)^{1/6}>0.2953$ (\Cref{prop:hand}). For all triples, the lengths of the steps and the angles
between consecutive steps are independent (\Cref{lem:angles}). Grouping the lengths into $97$ classes and the angles
into $64$ classes gives a non-negative matrix $P$ on $602{,}176$ states, indexed by two consecutive length classes and
one angle class, such that $q_m$ is at most a constant times the sum of the entries of $P^{m-2}$ applied to a fixed
positive vector (\Cref{prop:transfer}). A positive vector $v$ with $Pv\le\rho v$ coordinatewise then gives
$q_m\le C\rho^{m-2}$ for a constant $C$, and a computation in exact integer arithmetic produces such a
vector with $\rho=N/10^{12}$, $N=819{,}869{,}595{,}611$.

\begin{theorem}\label{thm:main}
For $d=2$, $c(0+)\ge b_2\,(10^{12}/N)^{1/2}=0.324170\dots$, so $c(0+)>0.32417$.
\end{theorem}

\Cref{sec:transfer} proves that the inequalities checked by the programs imply the theorem, and
\Cref{sec:computation} describes the computation and its independent verification.

\section{The first moment}\label{sec:first}

\begin{proof}[Proof of \Cref{thm:first}]
By the multivariate Mecke equation \cite[Theorem~4.4]{LP},
\[
\E Z_m(a)=\int_{(\R^d)^m}\mathbf 1\Bigl\{\sum_{i=1}^m|x_i-x_{i-1}|\le am\Bigr\}\,dx_1\cdots dx_m .
\]
The change of variables $u_i=x_i-x_{i-1}$ has Jacobian one, and in polar coordinates $du_i=d v_d r_i^{d-1}dr_i\,
d\sigma_i$ with $\sigma_i$ the normalized surface measure. Hence
\[
\E Z_m(a)=(dv_d)^m\int_{r_i\ge 0,\ \sum r_i\le am}\ \prod_{i=1}^m r_i^{d-1}\,dr
=(dv_d)^m\frac{\Gamma(d)^m (am)^{dm}}{\Gamma(dm+1)}=\frac{(v_d\,d!)^m (am)^{dm}}{(dm)!},
\]
by the Dirichlet integral. Stirling's formula gives
$\frac1m\log\E Z_m(a)\to\log\bigl(v_d\,d!\,(ae/d)^d\bigr)$, which is negative for $a<b_d$ and positive for $a>b_d$.
If $c(0+)<b_d$, choose $c(0+)<a<b_d$; then $\Prob(T_m\le am)\to 1$ by (e), while
$\Prob(T_m\le am)\le\E Z_m(a)\to 0$. So $c(0+)\ge b_d$.
\end{proof}

The same bound follows from the exponential tilt used in \cite{AK}: for $\lambda>0$,
\begin{equation}\label{eq:tilt}
\E Z_m(a)\le e^{\lambda am}\Bigl(\int_{\R^d}e^{-\lambda|u|}du\Bigr)^m=\bigl(e^{\lambda a}\,v_d\,d!\,\lambda^{-d}\bigr)^m,
\end{equation}
minimized at $\lambda=d/a$, where the base is $v_d\,d!\,(ae/d)^d$. We use the tilt below because it factorizes over
steps.

\begin{proof}[Proof of \Cref{cor:asym}]
As $d\to\infty$, Stirling's formula gives $\log v_d=\frac d2\log\frac{2\pi e}{d}-\frac12\log(\pi d)+o(1)$ and
$\log d!=d\log\frac de+\frac12\log(2\pi d)+o(1)$, so, with absolute implied constants, $(v_d\,d!)^{1/d}=\sqrt{2\pi d/e}\,(1+O(\log d/d))$. Hence
$b_d=\sqrt{d/(2\pi e)}\,(1+O(\log d/d))$, and with \Cref{thm:first}, $c(0+)\le\beta_d$ and \cite{Rhee} the limit
follows. For the second statement, $(v_d\,\Gamma(d+1))^{-1/(d-1)}=(v_d\,d!)^{-1/d}(1+O(\log d/d))$, so
$e(d-1)(v_d\,\Gamma(d+1))^{-1/(d-1)}=e^2 b_d(1+O(\log d/d))$, whose ratio to $\sqrt d$ tends to
$e^2(2\pi e)^{-1/2}>(2\pi e)^{-1/2}=\lim\beta_d/\sqrt d$.
\end{proof}

At $d=10^3$, $b_d/\sqrt{d/(2\pi e)}=0.99965\dots$ and
$e(d-1)(v_d\,\Gamma(d+1))^{-1/(d-1)}/\sqrt{d/(2\pi e)}=7.350\dots$ (computed by \texttt{first\_moment.py}).

\section{Shortest paths and the exchange of two points}\label{sec:exchange}

From now on $d=2$. A shortest path exists almost surely: if some path through $m$ points has length $B$, every path
of length at most $B$ uses points in the ball of radius $B$ about the origin, of which there are finitely many.

\begin{lemma}\label{lem:swap}
Let $(0=x_0,x_1,\dots,x_m)$ be a shortest path and $u_i=x_i-x_{i-1}$. For $1\le i\le m-2$,
\begin{equation}\label{eq:swap}
|u_i|+|u_{i+2}|\le|u_i+u_{i+1}|+|u_{i+1}+u_{i+2}|.
\end{equation}
\end{lemma}

\begin{proof}
Replacing $(x_{i-1},x_i,x_{i+1},x_{i+2})$ by $(x_{i-1},x_{i+1},x_i,x_{i+2})$ keeps the vertex set, both endpoints and
all other edges, and changes the length of this stretch from $|u_i|+|u_{i+1}|+|u_{i+2}|$ to
$|u_i+u_{i+1}|+|u_{i+1}|+|u_{i+1}+u_{i+2}|$.
\end{proof}

Call a triple of vectors $(u_1,u_2,u_3)$ \emph{bad} if $|u_1+u_2|+|u_2+u_3|<|u_1|+|u_3|$. Let $\mu_\lambda$ be the probability measure
on $\R^2$ with density $e^{-\lambda|u|}\lambda^2/(2\pi)$: its radius has the Gamma$(2,\lambda)$ law, its direction
is uniform, and the two are independent. Since badness is invariant under scaling, the probability that a triple of
independent $\mu_\lambda$ vectors is bad does not depend on $\lambda$.

\begin{proposition}\label{prop:block}
Let $U_1,U_2,\dots$ be independent with law $\mu_1$, and let $q_m$ be the probability that none of the triples
$(U_i,U_{i+1},U_{i+2})$, $1\le i\le m-2$, is bad. Then for all $a,\lambda>0$,
\[
\Prob(T_m\le am)\le\bigl(e^{\lambda a}\,2\pi\lambda^{-2}\bigr)^m q_m ,
\]
and if $\limsup_m q_m^{1/m}\le\rho$, then $c(0+)\ge b_2\,\rho^{-1/2}$.
\end{proposition}

\begin{proof}
On $\{T_m\le am\}$ there is a shortest path of length at most $am$, and by \Cref{lem:swap} none of its step triples is
bad. Let $Z^*_m(a)$ count the ordered $m$-tuples of distinct points of $\Pi$ whose path from the origin has length at
most $am$ and no bad triple of consecutive steps. By the Mecke formula and the change to increments, and since
$\mathbf 1\{\sum|u_i|\le am\}\le e^{\lambda(am-\sum|u_i|)}$,
\[
\Prob(T_m\le am)\le\E Z^*_m(a)\le e^{\lambda am}\int\prod_{i=1}^m e^{-\lambda|u_i|}\,
\mathbf 1\{\text{no bad triple}\}\,du=e^{\lambda am}(2\pi\lambda^{-2})^m q_m ,
\]
where the last step uses the scale invariance of badness. With $\lambda=2/a$ the base is $\pi e^2a^2/2$. If
$a<b_2\rho^{-1/2}$, then $(\pi e^2a^2/2)\rho<1$, the right side tends to $0$, and (e) gives $c(0+)\ge a$.
\end{proof}

Badness of disjoint triples involves disjoint sets of the $U_i$, so $q_m\le(1-\delta)^{\lfloor m/3\rfloor}$ whenever
$\delta$ is a lower bound for the probability that one triple is bad, and then $c(0+)\ge b_2(1-\delta)^{-1/6}$.

\begin{proposition}\label{prop:hand}
$c(0+)\ge b_2\,(27/26)^{1/6}>0.2953$.
\end{proposition}

\begin{proof}
Let $u_1,u_2,u_3$ be independent with law $\mu_1$. Let $F$ be the event that $|u_2|<\min(|u_1|,|u_3|)$ and that $u_1$ and $u_3$
both make an angle less than $\pi/3$ with $-u_2$. The radii and the directions are independent, the first condition
has probability $1/3$, and given the direction of $u_2$ the two angular conditions are independent with probability
$1/3$ each, so $\Prob(F)=1/27$. On $F$, $u_1\cdot u_2<-|u_1||u_2|/2$, so
\[
|u_1+u_2|^2<|u_1|^2+|u_2|^2-|u_1||u_2|<|u_1|^2,
\]
and likewise $|u_2+u_3|<|u_3|$; the triple is bad. Take $\delta=1/27$.
\end{proof}

\section{The transfer operator}\label{sec:transfer}

We now bound $q_m$ using all overlapping triples. Write $R_i=|U_i|$ and let $\varphi_i\in[0,2\pi)$ be the angle from
the direction of $U_i$ to that of $U_{i+1}$, so that $|U_i+U_{i+1}|^2=R_i^2+R_{i+1}^2+2R_iR_{i+1}\cos\varphi_i$.

\begin{lemma}\label{lem:angles}
The random variables $R_1,R_2,\dots,\varphi_1,\varphi_2,\dots$ are independent. Each $R_i$ has the Gamma$(2,1)$
law and each $\varphi_i$ is uniform on $[0,2\pi)$.
\end{lemma}

\begin{proof}
The directions $\theta_i\in[0,2\pi)$ of the vectors $U_i$ are independent, uniform, and independent of the radii. The map
$(\theta_1,\dots,\theta_k)\mapsto(\theta_1,\theta_2-\theta_1,\dots,\theta_k-\theta_{k-1})$ (mod $2\pi$) is a group
automorphism of the torus and preserves Haar measure.
\end{proof}

The triple $(U_i,U_{i+1},U_{i+2})$ is bad exactly when
\begin{equation}\label{eq:bad}
G(R_i,R_{i+1},\cos\varphi_i)+G(R_{i+2},R_{i+1},\cos\varphi_{i+1})<0,\qquad G(r,z,h)=\sqrt{r^2+z^2+2rzh}-r .
\end{equation}
For $|h|\le 1$ the function $G$ is non-increasing in $r\ge 0$ (it is $|(r,0)+z(h,\sqrt{1-h^2})|-r$, and the norm of
$(r,0)+y$ grows at most at rate one in $r$), non-decreasing in $h$, and convex in $z$.

\emph{Categories.} Fix integers $n,s,A\ge 1$, $Q$ and $D$. Radius bin $j\in\{0,\dots,n-1\}$ is $I_j=[j/s,(j+1)/s]$.
Choose integers $w_j\ge 0$ with $w_j\le Q\,\Prob(R\in I_j)$ and put $w_n=Q-\sum_{j<n}w_j>0$. Angle bin
$t\in\{0,\dots,A-1\}$ is $J_t=[2\pi t/A,2\pi(t+1)/A)$, and $C_t$ is an integer with $D\max_{\varphi\in J_t}
\cos\varphi\le C_t\le D$. On an enlarged probability space we attach to each $R_i$ a category $K_i\in\{0,\dots,n\}$, independent
across $i$, with $\Prob(K_i=j)=w_j/Q$ and $R_i\in I_j$ whenever $K_i=j<n$; this is possible because $w_j/Q\le
\Prob(R\in I_j)$ (retain a fraction of the mass of each bin, using an independent uniform variable, and assign the
rest to the category $n$). Let $t_i$ be the bin of $\varphi_i$. By \Cref{lem:angles}, the $K_i$ and the $t_i$ are
independent, with $\Prob(t_i=t)=1/A$.

\emph{The swap test.} For ordinary categories $a,b,c<n$ and angle bins $t,u$, call the cell $(a,b,c,t,u)$
\emph{forbidden} if, for both $z\in\{b,b+1\}$, the integers
\[
X=D(a^2+z^2)+2azC_t,\qquad Y=D(c^2+z^2)+2czC_u,\qquad L^2=D(a+c)^2
\]
satisfy $\sqrt X+\sqrt Y<L$. Put $H(a,b,c,t,u)=0$ on forbidden cells and $H=1$ otherwise, in particular
whenever one of $a,b,c$ equals $n$.

\begin{lemma}\label{lem:cell}
If $(K_i,K_{i+1},K_{i+2},t_i,t_{i+1})$ is a forbidden cell, then $(U_i,U_{i+1},U_{i+2})$ is bad.
\end{lemma}

\begin{proof}
Both sides of \eqref{eq:bad} are homogeneous, so we may multiply all radii by $s$; then $sR_i\in[a,a+1]$,
$sR_{i+1}\in[b,b+1]$ and $sR_{i+2}\in[c,c+1]$. By the monotonicity of $G$, the left side of \eqref{eq:bad} is at most
$G(a,z,C_t/D)+G(c,z,C_u/D)$ with $z=sR_{i+1}$, and this is convex in $z$, so it is at most its larger value at
$z\in\{b,b+1\}$. At such an endpoint, multiplying by $\sqrt D$ shows that $G(a,z,C_t/D)+G(c,z,C_u/D)<0$ is the
inequality $\sqrt X+\sqrt Y<L$.
\end{proof}

\emph{The operator.} Let $\mathcal S=\{0,\dots,n\}^2\times\{0,\dots,A-1\}$ and define $M$ on functions on
$\mathcal S$ by
\[
(Mv)(b,c,u)=w_c\sum_{a=0}^{n}\sum_{t=0}^{A-1}H(a,b,c,t,u)\,v(a,b,t).
\]
Put $V_2(a,b,t)=w_aw_b$ and $P=M/(QA)$.

\begin{proposition}\label{prop:transfer}
For $m\ge 2$, $q_m\le Q^{-2}A^{-1}\sum_{\mathcal S}P^{m-2}V_2$. If $v>0$ on $\mathcal S$ and $Pv\le\rho v$
coordinatewise, then $q_m\le C\rho^{m-2}$ for a constant $C$ independent of $m$, and $c(0+)\ge b_2\,\rho^{-1/2}$.
\end{proposition}

\begin{proof}
Let $E_m$ be the event that no cell $(K_i,K_{i+1},K_{i+2},t_i,t_{i+1})$, $1\le i\le m-2$, is forbidden. By
\Cref{lem:cell}, no bad triple implies $E_m$, so $q_m\le\Prob(E_m)$. Since the $K_i$ and $t_i$ are independent with
the stated laws, the sequence $(K_i,K_{i+1},t_i)$ is a Markov chain, and induction on $m$ shows that $Q^mA^{m-1}
\Prob(E_m,K_{m-1}=a,K_m=b,t_{m-1}=t)=(M^{m-2}V_2)(a,b,t)$; the initial vector accounts for two radius categories and
one angle bin, and each step adds one of each. This is the first claim. If $Pv\le\rho v$ with $v>0$, choose $C_0$ with
$V_2\le C_0v$; since $P$ has non-negative entries, $P^{m-2}V_2\le C_0\rho^{m-2}v$, which gives the second claim with $C=C_0Q^{-2}A^{-1}\sum_{\mathcal S}v$, and
the third follows from \Cref{prop:block}.
\end{proof}

\section{The computation}\label{sec:computation}

We take $n=96$, $s=12$ (radii up to $8$), $A=64$, $Q=2^{32}$ and $D=2^{30}$, so $|\mathcal S|=97^2\cdot 64=602{,}176$.
All inputs are exact.

\emph{Radius weights.} The Gamma$(2,1)$ law has $\Prob(R>x)=S(x)=(1+x)e^{-x}$, and
$w_j=\lfloor Q\,(S^-(j/s)-S^+((j+1)/s))\rfloor$ for rational bounds $S^-\le S\le S^+$ obtained from the Taylor
polynomial of $e^x$ of degree $80$ with an explicit tail bound. Then $w_n=Q-\sum_{j<n}w_j>0$ holds the tail beyond
radius $8$ and all rounding losses.

\emph{Cosine bounds.} For $j=\min(t,A-1-t)$ the largest value of $\cos$ on $J_t$ is $\cos(2\pi j/A)$ with
$0\le 2\pi j/A<\pi$. With a rational $\pi^-<\pi$ from Machin's formula $\pi=16\arctan\frac15-4\arctan\frac1{239}$ and
alternating series, $\cos(2\pi j/A)\le\cos(2\pi^-j/A)$, which is bounded above by the cosine series through degree
$80$ plus the absolute value of the next term; $C_t$ is $D$ times this bound, rounded up. For every $t$, $|C_t|\le D$.

\emph{The certificate.} The program \texttt{cert.cpp} finds a positive integer vector $v$ by power iteration in
floating point, then computes $Mv$ exactly in $128$-bit integer arithmetic (all integers involved are below $2^{126}$) and
reports the least integer $N$ with
\[
10^{12}\,(Mv)_i\le N\,Q\,A\,v_i\qquad\text{for all }i\in\mathcal S .
\]
The result is $N=819{,}869{,}595{,}611$, with $29{,}785{,}462\le v_i\le 2^{40}+1$. By
\Cref{prop:transfer} with $\rho=N/10^{12}$,
\[
c(0+)\ge\frac{\sqrt{2/\pi}}{e}\sqrt{\frac{10^{12}}{N}}=0.3241702746\dots,
\]
and the inequality $b_2\,\rho^{-1/2}>0.32417$ is equivalent to $2>0.32417^2\,\pi e^2\rho$, checked in rational arithmetic with
upper bounds for $\pi$ and $e$. The swap test is decided in integers as $\Delta>0$ and $\Delta^2>4XY$ with $\Delta=L^2-X-Y$,
which is equivalent to $\sqrt X+\sqrt Y<L$ because $(\sqrt X+\sqrt Y)^2=X+Y+2\sqrt{XY}$. The way $v$ was found plays no role: any positive $v$ for which the displayed
inequalities hold proves the bound.

\emph{Independent check.} The program \texttt{check.py} shares no code with \texttt{cert.cpp} and \texttt{grid.py}.
It re-verifies every weight and cosine bound with different enclosures: $e^{-x}$ by its alternating series, $\pi$ by
Euler's formula $\pi=4\arctan\frac12+4\arctan\frac13$, and $\cos x=1-2\sin^2(x/2)$ with a lower bound for $\sin$. It
then recomputes $Mv$ from $v$ without the threshold structure used in \texttt{cert.cpp}. The swap test is evaluated in
floating point, and every comparison within $10^{-3}$ of a tie would be redone in integers (in the run reported here there were none), in the form $L^2>X+Y$ and
$(L^2+Y-X)^2>4L^2Y$. Given $L^2>X+Y$, the second inequality is $L^2+Y-X>2L\sqrt Y$, that is, $X<(L-\sqrt Y)^2$, which
is again $\sqrt X+\sqrt Y<L$. The sums are matrix products on the $20$-bit halves of $v$, which are exact in double precision.
The result is compared with the saved $Mv$ entry by entry, and the $602{,}176$ inequalities and the final rational
inequality are checked with exact integers.

\section*{Acknowledgments and funding}

This work was done by multiple AI agents, Claude (Anthropic) and Codex (OpenAI), connected to Math
Brain~\cite{MathBrainRepo} and working together inside the Brain framework~\cite{BrainFramework}, a shared, persistent
memory in which the agents record claims, proofs, failed attempts and methods, retrieve one another's results,
coordinate who works on what, and audit each other's work. For more on the Brain and Math Brain, see the Math Brain
paper~\cite{MathBrainPaper}, the Math Brain repository~\cite{MathBrainRepo} and the Brain framework
repository~\cite{BrainFramework}.

The research presented in this paper was supported by the European Research Council (ERC) under the European
Union's Horizon Europe research and innovation programme (grant agreement No.\ 101041711). It was also supported by
the Simons Foundation (as part of the Collaboration on the Mathematical and Scientific Foundations of Deep Learning),
by Heights Labs, by Convex Nexus Capital, and by the Israel Science Foundation (grants 2258/19 and 4101/25).

\section*{Declaration of competing interest}

The author holds equity in Heights Labs and Convex Nexus Capital and has a management role at Convex Nexus Capital.

\section*{Declaration of generative AI and AI-assisted technologies}

The arguments of \Cref{sec:first,sec:exchange,sec:transfer} and the programs of \Cref{sec:computation} were developed
with AI systems (OpenAI Codex and Anthropic Claude agents working under the author's direction), which also drafted the
text. The problem, the limit theorem used throughout and the branching random walk comparison are due to Aldous and
Krikun \cite{AK}; the path count with an exponential weight is in the proof of \cite[Lemma~2.1]{GM} of Gou\'er\'e
and Marchand; the asymptotic upper bound is Rhee's theorem \cite{Rhee}. The author directed the work and is
responsible for the content of the paper and for decisions concerning its dissemination.

\section*{Data and code availability}

The programs \texttt{grid.py}, \texttt{cert.cpp}, \texttt{check.py} and \texttt{first\_moment.py}, the grid file and
the checksums and the vectors $v$ and $Mv$ of the certificate are in the folder \texttt{code} of the repository
\begin{center}\url{https://github.com/yspennstate/sparse-tsp-constant-lower-bound}.\end{center}

\appendix
\section{Notation}\label{app:notation}

\begin{center}
\small
\begin{tabular}{lll}
symbol & meaning & first use\\\hline
$\Pi$, $\xi_i$ & Poisson process of unit intensity in $\R^d$, its points & \Cref{sec:intro}\\
$\Prob$, $\E$, $\mathbf{1}$ & probability, expectation, indicator & \Cref{sec:intro}\\
$\Gamma$ & Euler's gamma function & \Cref{sec:intro}\\
$v_d$ & volume of the unit ball in $\R^d$ & \Cref{sec:intro}\\
$T_m$ & shortest path length from $0$ through $m$ points of $\Pi$ & \Cref{sec:intro}\\
$L_k(\delta)$ & shortest cycle through $\lceil\delta k\rceil$ of $k$ points & \Cref{sec:intro}\\
$c(\delta)$, $c(0+)$ & the constants of Aldous and Krikun & \Cref{sec:intro}\\
$\beta_d$ & the travelling salesman constant, $c(1)$ & \Cref{sec:intro}\\
$W(\ell)$ & most points of $\Pi$ on a path of length at most $\ell$ & \Cref{sec:intro}\\
$b_d$ & first moment bound $(d/e)(v_d\,d!)^{-1/d}$ & \Cref{sec:intro}\\
$Z_m(a)$, $Z^*_m(a)$ & numbers of short paths, with no bad triple for $Z^*_m$ & \Cref{sec:intro}, \Cref{sec:exchange}\\
$\mu_\lambda$ & law with density $e^{-\lambda|u|}\lambda^2/(2\pi)$ on $\R^2$ & \Cref{sec:exchange}\\
$q_m$ & probability of no bad triple among $m$ steps & \Cref{sec:intro}, \Cref{prop:block}\\
$U_i$ & independent steps with law $\mu_1$ & \Cref{prop:block}\\
$R_i$, $\varphi_i$ & step length, angle between consecutive steps & \Cref{sec:transfer}\\
$\theta_i$ & direction of the step $U_i$ & \Cref{sec:transfer}\\
$I_j$, $J_t$ & radius bin, angle bin & \Cref{sec:transfer}\\
$X$, $Y$, $L$ & integer squared lengths in the exchange test & \Cref{sec:transfer}\\
$\Delta$ & $L^2-X-Y$ & \Cref{sec:computation}\\
$S$ & tail $\Prob(R>x)=(1+x)e^{-x}$ of the step length & \Cref{sec:computation}\\
$G(r,z,h)$ & $\sqrt{r^2+z^2+2rzh}-r$ & \eqref{eq:bad}\\
$n,s,A,Q,D$ & grid parameters & \Cref{sec:transfer}\\
$w_j$, $C_t$ & integer radius weights, cosine bounds & \Cref{sec:transfer}\\
$K_i$, $t_i$ & radius category, angle bin & \Cref{sec:transfer}\\
$\mathcal S$, $V_2$ & state space, initial vector & \Cref{sec:transfer}\\
$H$, $M$, $P$ & cell indicator, transfer operator, $M/(QA)$ & \Cref{sec:intro}, \Cref{sec:transfer}\\
$E_m$ & no forbidden cell among $m$ steps & \Cref{prop:transfer}\\
$v$, $N$, $\rho$ & certificate vector, its numerator, $N/10^{12}$ & \Cref{sec:intro}, \Cref{sec:computation}
\end{tabular}
\end{center}
Letters used inside a single proof (such as $B$, $F$ and $u_1,u_2,u_3$) are defined there and are not listed.

\begin{thebibliography}{99}

\bibitem{Aapp}
D.~Aldous, \emph{Upper bounds for $c(0+)$}, appendix to \cite{AK}, November 2006.
\url{https://www.stat.berkeley.edu/~aldous/Papers/me110-app.pdf}

\bibitem{Anote}
D.~Aldous, correction to \cite[(3.6)]{AK}, note on the author's web page, accessed 19 November 2026.
\url{https://www.stat.berkeley.edu/~aldous/Papers/me110-update.pdf}

\bibitem{AK}
D.~Aldous and M.~Krikun, \emph{Percolating paths through random points}, ALEA Lat.\ Am.\ J.\ Probab.\ Math.\ Stat.\
\textbf{1} (2006), 89--109. arXiv:math/0509492.

\bibitem{BHH}
J.~Beardwood, J.~H.~Halton and J.~M.~Hammersley, \emph{The shortest path through many points}, Proc.\ Cambridge
Philos.\ Soc.\ \textbf{55} (1959), 299--327. \doi{10.1017/S0305004100034095}.

\bibitem{CBDM}
J.~G.~Carlsson, M.~Behroozi, R.~Devulapalli and X.~Meng, \emph{Household-level economies of scale in transportation},
Oper.\ Res.\ \textbf{64} (2016), 1372--1387. \doi{10.1287/opre.2016.1533}.

\bibitem{Croes}
G.~A.~Croes, \emph{A method for solving traveling-salesman problems}, Oper.\ Res.\ \textbf{6} (1958), 791--812.
\doi{10.1287/opre.6.6.791}.

\bibitem{GG}
J.~Gaudio and C.~K.~Guan, \emph{An improved upper bound for the Euclidean TSP constant using band crossovers}, 2026.
arXiv:2602.11250.

\bibitem{GM}
J.-B.~Gou\'er\'e and R.~Marchand, \emph{Continuous first-passage percolation and continuous greedy paths model: linear
growth}, Ann.\ Appl.\ Probab.\ \textbf{18} (2008), 2300--2319. \doi{10.1214/08-AAP523}.

\bibitem{LP}
G.~Last and M.~Penrose, \emph{Lectures on the Poisson Process}, Cambridge University Press, Cambridge, 2017.

\bibitem{Rhee}
W.~T.~Rhee, \emph{On the travelling salesperson problem in many dimensions}, Random Structures Algorithms \textbf{3}
(1992), 227--233. \doi{10.1002/rsa.3240030302}.

\bibitem{Verges}
J.~Verges, \emph{Concentration inequalities and large deviations for continuous greedy animals and paths}, 2025.
arXiv:2503.01484.

\bibitem{MathBrainRepo}
Y.~Shmalo, \emph{Math Brain: public snapshot, code and tests}, 2026.
\url{https://github.com/yspennstate/math-brain-public}

\bibitem{MathBrainPaper}
Y.~Shmalo, \emph{Math Brain: shared, auditable memory for multi-agent mathematical research}, October 2026.
\url{https://github.com/yspennstate/math-brain-public/tree/main/papers/math-brain-2026-10-24}

\bibitem{BrainFramework}
Y.~Shmalo, \emph{Brain Framework}, public software repository, 2026.
\url{https://github.com/yspennstate/brain-framework-public}

\end{thebibliography}

\end{document}
